\section{Worked problems: inference is paid per token}

\subsection{When service overtakes training}
\textbf{Problem.} A model with $N$ active parameters is trained on
$D = 30$~trillion tokens. Using $6ND$ FLOPs for training and $2N$ per processed
token in service, how many tokens must it process in service before inference
has spent as much compute as training? Why does the answer not depend on $N$,
and how does a prompt cache that serves a fraction $h$ of input tokens without
computing them change it?

\textbf{Solution.} Service matches training when $2N \cdot T = 6ND$, so
$T = 3D = 90$~trillion processed tokens. $N$ cancels because both costs are
proportional to it: every parameter is used twice per token in service and six
times per token in training. A prompt cache does not change the break-even in
\emph{computed} tokens, but it changes how many tokens users must send to reach
it: if a fraction $h$ of input tokens is served from the cache, then
$T_{\text{served}} = 3D / (1 - h \cdot f_{\text{in}})$, where $f_{\text{in}}$ is
the input share of all tokens. With DeepSeek's reported day --- 608~billion
input tokens, 56\% cached, and 168~billion output tokens ---
$f_{\text{in}} = 608/776 = 0.78$ and $1 - 0.56 \times 0.78 = 0.56$: users must
send about 1.8 times as many tokens before the service has repeated the
training run's arithmetic.

\subsection{What doubling the bandwidth would do}
\textbf{Problem.} On the M1, Llama~3.2~3B decoded one conversation at
20.08 tokens per second and thirty-two at 175.87 in total. Compute the step time
in each case. Then predict both numbers on a machine with twice the memory
bandwidth and the same arithmetic units, stating which assumption each
prediction rests on.

\textbf{Solution.} Step time is the batch divided by the aggregate rate:
$1/20.08 = 49.8$~ms at $B = 1$ and $32/175.87 = 182$~ms at $B = 32$. At
$B = 1$ the step is almost entirely the time to stream 2.01~GB of weights
(\Chref{roofline}), so twice the bandwidth predicts about 25~ms and
40~tokens per second --- assuming the kernel keeps its share of the machine's
bandwidth and the fixed per-step overheads are small. At $B = 32$ the step
already spends about 132~ms more than the weight stream, on arithmetic and on
reading 32 caches. Doubling bandwidth shortens only the memory part: if those
132~ms are mostly arithmetic, the step falls to about $182 - 25 = 157$~ms, or
204 tokens per second in total --- 16\% more rather than twice. The faster memory helps the single user most and the busy server
least, which is the roofline's prediction in miniature.

\subsection{A ledger for a familiar cache}
\textbf{Problem.} Write the constraint ledger for caching rendered web pages
in a web server's memory, using the book's four resources and two outcomes.

\textbf{Solution rubric.} A good ledger has one row per resource the cache
touches and a size or condition for each. \emph{Compute} and \emph{latency}:
buys --- a hit skips rendering; the saving is the render time multiplied by the
hit rate. \emph{Memory capacity}: spends --- every cached page occupies
memory, which bounds how many pages can be kept and competes with other uses.
\emph{Memory bandwidth}: roughly neutral --- a hit still copies the page out.
\emph{Correctness}: risks --- a stale page after the source changes, or one
user's personalized page served to another; the ledger should name the
invalidation rule. \emph{Cost}: buys, if the hit rate is high enough that the
memory costs less than the rendering it avoids. Every item has an inference
analogue: the prefix cache of \Chref{prefix-caching} buys compute and latency,
spends capacity, and risks leaking one tenant's prompt to another
(\Chref{production}).
